\documentclass[11pt]{article}
\usepackage[margin=1in]{geometry}
\usepackage{amsmath,amssymb,amsthm,booktabs,microtype,hyperref}
\hypersetup{hidelinks,pdfauthor={T},pdftitle={Separating Consequence and Tone in Preference Feedback}}
\newtheorem{proposition}{Proposition}
\title{Separating Consequence and Tone in Preference Feedback}
\author{T}
\date{}
\begin{document}
\maketitle

\begin{abstract}
Human feedback can reward both relief of a
problem and the feeling of having been heard. We specify a factorial
feedback design that estimates both contributions and exposes a ranking
reversal. Two structural models with identical observed training examples
have different causal effects. The result motivates interventional
examples during reward-model development and an explicit outcome holdout.
Three fitted affine scorers illustrate the role of the training target.
An outcome target chose benefit in all 36 synthetic holdout pairs.
The pooled rating target chose benefit in twelve and tied in one.
Two opposite-corner comparisons identify average effects but can
leave the effect of helping ambiguous at a fixed tone. We derive
sharp bounds for that effect from the two comparisons. When the
comparisons are reported as intervals, a graph gives exact effect
bounds and detects incompatible panels.
\end{abstract}

\section*{Summary and scope}
\paragraph{Question.} What must a preference study vary to separate
genuine help from a tone cue?
\paragraph{Finding.} Opposite-corner comparisons identify average
effects. The missing factorial interaction can still reverse the
consequence effect at one tone. Four independently defined cells
resolve that ambiguity in the rating table. An interval-valued panel
can sometimes certify the sign using a single extra comparison.
\paragraph{Evidence and scope.} The bounds are exact and
checked by enumeration. The training results use synthetic consequence
labels and supplied features. They characterize those learning targets
and leave the validity of real outcome labels to separate adjudication.

\section{What a high rating buys}
A warm dismissal can receive a higher immediate rating than a plain
answer that solves the problem. Comfort can itself be valuable. The
development risk is teaching a tone cue to substitute for the sought
help. Constitutional AI demonstrates a training pipeline using
critiques, revision and preference feedback \cite{bai}. Reward-gaming
analysis makes the general proxy problem precise \cite{skalse}. Here
the intervention changes the example design upstream of that proxy.

Let $h$ be a declared helpful consequence and $t$ a tone treatment,
both binary. In a synthetic crossed panel take rating $R(h,t)=2h+3t$.
At fixed tone, helping raises the score by two in both cells. Yet
$R(0,1)=3>2=R(1,0)$. A scalar fitted perfectly to this rating can
reproduce that reversal. Crossed grading makes the tradeoff visible
even when the fit looks excellent.

\section{What two paired comparisons can establish}
Suppose the assessor keeps only two comparisons. In the first,
the useful response is warm and the unhelpful one is plain.
In the second, those tone assignments are reversed. Write
$R_{ht}\in[0,1]$ for the normalized mean rating, with $h=+$
denoting the useful response and $t=+$ the warm tone. The
observed contrasts are
\[
 A=R_{++}-R_{--},\qquad B=R_{+-}-R_{-+}.
\]
Their half-sum is the average effect of the useful response
across tones, and their half-difference is the average tone
effect across consequences. Neither contrast gives a direct
comparison of consequences at a fixed tone.

\begin{proposition}
Let $W=2-|A|-|B|$. With only these two contrasts and the unit
rating range, the sharp interval for the factorial interaction
$C=R_{++}+R_{--}-R_{+-}-R_{-+}$ is $[-W,W]$.
The two fixed-tone consequence effects are
$(A+B+C)/2$ and $(A+B-C)/2$. Each has sharp interval
\[
 \left[\frac{A+B-W}{2},\frac{A+B+W}{2}\right].
\]
\end{proposition}
\begin{proof}
Two numbers in $[0,1]$ whose difference is $A$ can have
any sum in $[|A|,2-|A|]$. The other pair can independently
have any sum in $[|B|,2-|B|]$. Taking the difference of
these sums gives the interaction interval, with both endpoints
attainable. Expanding each fixed-tone contrast in terms of
$A,B,C$ yields the claimed effects and endpoints.
\end{proof}

For example, take $(A,B)=(0.6,0.2)$. The average consequence
effect is $0.4$, while the fixed-tone effect can be $-0.2$.
In order $(++,+-,-+,--)$, the table $(0.6,1,0.8,0)$ gives
effects $-0.2$ and $1$. The table $(0.8,0.6,0.4,0.2)$ gives
$0.4$ at both tones with the same observed contrasts. Measuring
one fixed-tone comparison in the same rating population identifies
the missing interaction. If the claims concern actual user benefit,
the outcome check still needs independent validation.

Across 17,667 four-cell tables on three rational grids, exact
enumeration recovered the predicted extrema for every contrast
pair. This finite check complements the continuous proof and
does not estimate how often the ambiguity occurs in human feedback.

\section{Combining comparisons reported as intervals}
Different reviewers may estimate different contrasts. Suppose each
reported interval bounds the difference of two cells of one declared
reference population. The intervals must already allow for sampling
uncertainty and any population differences the audit permits. We
ask which rating tables remain consistent with all of them.

Include a vertex for each cell and an anchor vertex with score zero.
For a bound $\ell\le R_i-R_j\le u$, add an edge $j\to i$ of
weight $u$ and an edge $i\to j$ of weight $-\ell$. Connect the
anchor to each cell and back with weights one and zero, respectively.
Write $d(i,j)$ for the shortest path length from $i$ to $j$.

\begin{proposition}
A common rating table exists exactly when the graph has no negative
cycle. In that case the sharp interval for $R_a-R_b$ is
$[-d(a,b),d(b,a)]$.
\end{proposition}
\begin{proof}
Each edge bounds a difference of scores. Adding the constraints
around a cycle makes its weight nonnegative. Conversely, if there
is no negative cycle, distances from any one vertex are finite
because all cells connect to the anchor. The triangle inequality
makes those distances a feasible set of scores after subtracting
the anchor distance. Adding the inequalities along any path gives
the stated interval. Using distances from $b$ attains the upper
endpoint, and using distances from $a$ attains the lower endpoint.
The feasible tables form a convex set, so every intermediate
value is also attained.
\end{proof}

For example, suppose the diagonals give $A\in[0.55,0.65]$ and
$B\in[0.15,0.25]$. The warm-tone consequence effect could be as
low as $-0.30$ or as high as $1$. A third comparison that puts
the plain-tone effect in $[0.50,0.60]$ narrows the warm-tone
interval to $[0.10,0.40]$. The sign is now determined even though
none of the three comparisons is exact. An audit can use this test
to decide whether paying for a fourth comparison could change the
sign conclusion under the declared intervals.

There is also a direct inconsistency check. With diagonals $0.60$
and $0.20$, the two fixed-tone effects must sum to $0.80$.
Intervals $[0.30,0.50]$ and $[0.70,0.90]$ each fit the rating
scale but have a combined lower bound of $1$. A cycle detects the
conflict. The diagnosis is incompatibility under the common-table
assumption. It cannot identify which panel, uncertainty interval,
or population assumption failed.

An exhaustive oracle classified 256 finite panels, of which 66
were incompatible. On the 190 feasible panels it confirmed all
4,560 lower and upper endpoints across ordered cell pairs. With
integer-scaled constraints the distance witnesses are integer-scaled
too, so the grid oracle covers the continuous extrema in these cases.

\section{An observational obstruction}
\begin{proposition}
There are two causal models with the same observed distribution of
action and outcome and different average effects of action.
\end{proposition}
\begin{proof}
Draw $X$ uniformly from $\{0,1\}$ and set $A=X$. In model I let $Y=A$.
In model II let $Y=X$. Both observational laws put half their mass on
each of $(0,0)$ and $(1,1)$. Replacing the action equation by $A=a$
sets $Y=a$ in I and leaves its mean $1/2$ in II. The average effects
are one and zero, respectively. This uses the structural intervention
semantics described by Pearl \cite{pearl}.
\end{proof}

\section{Development protocol}
Build four matched alternatives per scenario wherever the variation is
legitimate. Hold intent, privacy and permission constant. Independently
vary tone and the concrete effect, such as whether a correction
actually fixes the error. Record the beneficiary's stated goal and an
outcome check separately from satisfaction. A warm explanation can
win within the set that meets the task and respects the beneficiary.
In cases where tone changes uptake, keep the interaction term
$R(1,1)-R(1,0)-R(0,1)+R(0,0)$ rather than forcing separability.

Compare an ordinary preference reward against a reward with separately
supervised outcome and tone heads. Freeze scenario families and check
whether the latter reduces warm-but-unhelpful choices when the two
signals conflict. Include neutral-but-helpful and harmful-but-popular
alternatives, with independent outcome adjudication. Train only in
settings where synthetic or consented outcomes can be checked. Wrong
outcome labels remain dangerous under either design.

\section{Audit and next measurement}
The local check enumerated all four rating cells, confirmed contrasts
$(2,2)$, and evaluated both intervention distributions. The synthetic
effect magnitudes are a counterexample to identification from passive
ratings.

A separate training job fitted three affine scorers on the same four
balanced feature rows from zero initialization. Each used eight exact
gradient steps of size $1/2$. The feature columns are orthogonal and
have unit mean square. A closed-form least-squares oracle therefore
checks each fitted vector as $255/256$ times its optimum. The targets
were pooled rating, direct helpful outcome, and deliberately reversed
outcome. The 36 held-out pairs crossed new strength values for both
features. Each compares a helpful plain action with a warm unhelpful one.
\begin{center}
\begin{tabular}{lrrr}
\toprule
Training target & Chose help & Tied & Chose warm unhelpful\\
\midrule
Pooled rating & 12 & 1 & 23\\
Direct outcome & 36 & 0 & 0\\
Reversed outcome & 0 & 0 & 36\\
\bottomrule
\end{tabular}
\end{center}
The reversed target is the decisive source-validity check. The
synthetic model learns the goal developers specify, including a
wrong one. Here legitimate and proxy features are supplied by the
experimenter. A future model trial needs outcome fidelity and realistic
crossing. A manipulation that changes the apparent tone and what was
actually offered must be treated as a different intervention.
\section{Matched stochastic training study}
Eight train and holdout seed pairs tested whether the earlier result
survives noisy supervision. We fitted logistic scores from 320
balanced-in-expectation feature draws per seed. The pooled labels
were Bernoulli with logit \(0.9h+2t\), and direct labels had logit
\(2h\). The two arms used the same inputs, uniforms, optimizer,
and evaluation pairs. Each new pair compared a helpful colder choice
with a warm unhelpful choice at unseen feature strengths.

The direct outcome arm chose help on 3,200 of 3,200 pairs, compared
with 58 for pooled feedback. The paired gain ranged from 94.75 to
100 percentage points across the eight seeds. This extends the
deterministic example to noisy labels within the supplied feature
map. The latent target remains an experimenter definition.

\paragraph{Review status.} Independent manuscript review and a language-model study remain open.

\begin{thebibliography}{9}
\bibitem{bai} Yuntao Bai et al. \emph{Constitutional AI: Harmlessness from AI Feedback}.
\url{https://arxiv.org/pdf/2212.08073}.
\bibitem{skalse} Joar Skalse, Nikolaus H. R. Howe, Dmitrii Krasheninnikov,
and David Krueger. \emph{Defining and Characterizing Reward Hacking}.
\url{https://arxiv.org/pdf/2209.13085}.
\bibitem{pearl} J. Pearl. Causal inference in statistics: An overview.
\url{https://ftp.cs.ucla.edu/pub/stat_ser/r350.pdf}.
\end{thebibliography}
\end{document}
